\chapter{Le moteur de cotation}\label{ch:hf:the-quoting-engine}

Un teneur de marché modifie son prix plusieurs dizaines de milliers de fois par jour sur un instrument actif. Chaque modification est une annulation suivie d'un nouvel ordre, ou
une modification~; chacune peut croiser une exécution déjà en route, et la plateforme les comptabilise toutes. Le composant qui ramène «~où je veux être~» au
plus petit nombre de messages pour y parvenir est le moteur de cotation. Dans le marché simulé de ce chapitre, un cotateur à un seul niveau envoie 758 messages par heure pour 250
exécutions. Une échelle à trois niveaux en envoie 980 et gagne davantage, car ses ordres sont déjà dans la file lorsque le prix bouge. Une hystérésis qui ressemble à
de l'économie, consistant à conserver un ordre à un tick de sa cible, rend le mark-out négatif.

\section{Des cotations cibles aux ordres}

Les politiques des chapitres 3 à 9 produisent des cibles~: de chaque côté, les prix et les tailles que le teneur de marché souhaite afficher. Le moteur gère les ordres~:
lesquels sont au repos, lesquels sont en route, lesquels la plateforme n'a pas encore accusés de réception, lesquels elle a été priée d'annuler. Il compare la cible avec
les ordres et envoie la différence. Ses règles (\cref{lst:hf:the-quoting-engine:diff})~:

\begin{method}[Comparer une cible aux ordres en cours]\label{met:hf:the-quoting-engine:diff}
De chaque côté, parcourir les ordres en cours par priorité de prix. Un ordre à un prix cible couvre celui-ci~: s'il porte plus que la cible, le réduire~;
s'il porte moins, ajouter un nouvel ordre pour la différence, sans jamais l'augmenter. Un ordre à aucun prix cible est annulé, sauf s'il se trouve dans une hystérésis
d'une cible qu'il peut couvrir. Toute cible restée non couverte donne lieu à un nouvel ordre. Laisser de côté les ordres en attente d'accusé de réception. Envoyer d'abord les annulations, puis
les modifications, puis les nouveaux ordres.
\end{method}

L'ordre compte. Les annulations partent en premier, car une annulation protège le capital alors qu'un nouvel ordre l'engage~; lorsque les messages doivent être rationnés, on envoie
d'abord ceux qui retirent du risque. Les ordres en route sont laissés de côté, parce que le moteur ne sait pas encore s'ils sont dans le carnet~; agir sur eux deux fois est la façon dont
les moteurs produisent des ordres en double.

\begin{definition}[Échelle de cotations]\label{def:hf:the-quoting-engine:ladder}
Une \emph{échelle de cotations}\index{echelle@échelle de cotations} (quote ladder) est un ensemble d'ordres au repos à plusieurs prix consécutifs d'un même côté, chacun destiné à être exécuté~: le meilleur
niveau rapporte déjà le spread, et les niveaux qui le suivent occupent des places dans des files qui deviennent les meilleures lorsque le prix bouge.
\end{definition}

Dans un carnet à grand tick (chapitre~5), les niveaux plus profonds d'une échelle sont rarement exécutés à leur place~; leur valeur tient à la place qu'ils occupent. Lorsque le prix monte d'un tick,
un cotateur à un seul niveau annule son cours acheteur et rejoint la fin de la nouvelle file au meilleur cours~; le second niveau d'une échelle est déjà dans cette file.

\section{Annulation-remplacement, priorité et états en route}

\begin{definition}[Annulation-remplacement]\label{def:hf:the-quoting-engine:replace}
Une \emph{annulation-remplacement}\index{annulation-remplacement} (cancel-replace) modifie le prix ou la taille d'un ordre au repos en un seul message, la plateforme annulant l'ancien ordre
et en inscrivant un nouveau~; selon la règle habituelle des plateformes, l'ordre conserve sa place dans la file seulement si sa taille est réduite au même prix, et il passe en fin de file
pour tout autre changement.
\end{definition}

Le livre~10, chapitre~26, précise le message de remplacement du simulateur selon cette règle~: une diminution de taille au même prix est publiée comme
une annulation partielle et conserve la priorité~; tout autre changement la fait perdre. C'est pourquoi le moteur ne modifie qu'à la baisse et ajoute un ordre distinct pour augmenter la taille.

Un ordre a un cycle de vie que le moteur doit suivre (\cref{fig:hf:the-quoting-engine:states})~: envoyé, puis accusé de réception~; demande de modification ou d'annulation, puis nouvel accusé de réception. Une exécution peut arriver à tout
moment, y compris après que le moteur a décidé d'annuler l'ordre~: l'annulation est encore en transit, ou elle a atteint la plateforme juste après l'exécution, qui répond
«~annulation trop tardive~». Cette course entre annulation et exécution n'est pas une erreur~; à des taux de messages élevés, c'est une certitude, et le moteur la comptabilise. Dans les
simulations du chapitre, un cotateur à un seul niveau rencontre 19.5 courses de ce type par heure.

\begin{omfigure}
\begin{tikzpicture}[font=\footnotesize,>=stealth,st/.style={draw,rounded corners,fill=omThm!8,minimum width=2.3cm,minimum height=0.7cm}]
\node[st] (pn) at (0,0) {nouvel ordre};
\node[st] (lv) at (4.2,0) {actif};
\node[st] (pa) at (4.2,1.8) {modif. en attente};
\node[st] (pc) at (8.4,0) {annul. en attente};
\node[st,fill=gray!10] (gone) at (8.4,-1.8) {retiré};
\draw[->] (pn) -- node[above] {accusé} (lv);
\draw[->] (lv.north west) to[bend left=15] node[left] {réduction} (pa.south west);
\draw[->] (pa.south east) to[bend left=15] node[right] {accusé} (lv.north east);
\draw[->] (lv) -- node[above] {annulation} (pc);
\draw[->] (pc) -- node[right] {accusé} (gone);
\draw[->,dashed] (lv.south) |- node[pos=0.25,left] {exécution totale} (gone.west);
\draw[->,dashed,omDef] (pc.north east) to[out=40,in=320,looseness=5] node[right,align=left] {exécution~: la\\course annulation--exécution} (pc.south east);
\end{tikzpicture}

\omcaption{Les états du moteur pour un ordre. Les flèches pleines sont les actions du moteur et les accusés de réception de la plateforme~; les flèches en pointillés sont des exécutions,
qui peuvent arriver dans n'importe quel état~; une exécution reçue pendant qu'une annulation est en transit est comptée comme une course. Une augmentation de taille n'est jamais une modification~: c'est un
nouvel ordre au même prix. Source~: \texttt{firm.quoteengine}.}\label{fig:hf:the-quoting-engine:states}
\end{omfigure}

\section{Limiteurs, plafonds de messages et ratios ordres sur transactions}

\begin{definition}[Limiteur de messages]\label{def:hf:the-quoting-engine:throttle}
Un \emph{limiteur de messages}\index{limiteur de messages} (throttle) est une limite sur le débit des messages qu'un système de négociation envoie, imposée par la plateforme à chaque session ou par
la société à elle-même, ainsi que la politique qui définit ce qui part lorsque la limite est atteinte~: quels messages partent en premier, lesquels sont différés et lesquels sont abandonnés.
\end{definition}

Le limiteur du moteur est un seau à jetons (livre~13, chapitre~21, en construit un pour la passerelle)~: les jetons se rechargent au débit permis, jusqu'à une
profondeur de rafale~; chaque message en consomme un, et lorsqu'il n'en reste plus, les actions restantes sont abandonnées. Abandonner, et non mettre en file d'attente, est la bonne politique pour les cotations~: la
mise à jour suivante recalcule la différence à partir de la dernière cible, de sorte qu'une action périmée n'est jamais envoyée en retard.

Les plateformes limitent les messages de trois façons~: un débit maximal strict par session (un limiteur rejette l'excédent), une redevance sur les messages au-delà d'un ratio par rapport aux transactions, et le
\emph{ratio ordres sur transactions} réglementaire du livre~10, chapitre~24.

\begin{dated}{2026-09}{Limites de messages dans les règles et les barèmes de frais}\label{dat:hf:the-quoting-engine:limits}
En vertu du règlement délégué (UE) 2017/566 de la Commission, chaque plateforme de négociation calcule, pour chaque membre et chaque instrument, au moins à la fin de chaque
session, le ratio des ordres non exécutés sur les transactions en volume et en nombre (ordres divisés par transactions, moins un), en comptant chaque type d'ordre
selon ce que prescrit son annexe~; un membre qui dépasse le ratio maximal de la plateforme sur l'une ou l'autre mesure a dépassé ce seuil. Le programme d'efficacité des messages du CME Group,
tel que déposé auprès de la CFTC en 2019, pondère les messages par type en un score de messages, le divise par le volume négocié par la société dans un groupe de produits
pendant les heures de négociation régulières, et compare le ratio de volume obtenu aux indices de référence par groupe de produits fixés chaque trimestre.
\end{dated}

\section{Prévention de l'auto-appariement et stratégies multiples}

Une société qui fait tourner plusieurs stratégies sur un même instrument peut voir son propre cours acheteur croiser son propre cours vendeur~: une stratégie de cotation laissée au repos au cours vendeur et une
stratégie de momentum qui achète. La transaction est réelle dans les livres de la plateforme, mais elle ne déplace aucun risque, paie deux fois les frais et laisse une transaction fictive que la surveillance
remarquera (livre~9, chapitre~29). Les plateformes proposent la \emph{prévention de l'auto-appariement} (self-trade prevention) (livre~10, chapitre~26)~: les ordres marqués du même groupe
ne sont jamais appariés entre eux, et la plateforme annule le plus récent, le plus ancien ou les deux, selon le choix de la société. La part du moteur dans ce problème est de
connaître tous les ordres au repos de chaque stratégie et acheminer les croisements internes vers la logique de compensation du chapitre~11 avant qu'ils n'atteignent la plateforme.

\section{Échelles de cotations et la ligne à ne pas franchir}

Une échelle d'ordres d'un même côté, annulés et déplacés au gré du marché, ressemble de l'extérieur au layering du livre~9, chapitre
29~: des ordres à plusieurs prix qui sont rarement exécutés. La différence tient à l'intention. La loi en vertu de laquelle le septième circuit a confirmé la condamnation de Coscia en
2017 définit le spoofing comme «~offrir ou demander avec l'intention d'annuler l'offre ou la demande avant exécution~». La cour a rejeté son argument selon lequel
la définition était vague, puisque les traders à haute fréquence annulent 98~\% de leurs ordres. Annuler la plupart des ordres est normal~; les placer sans jamais vouloir
qu'ils soient exécutés ne l'est pas. L'échelle d'un teneur de marché est défendable lorsque chaque niveau est fixé en prix et en taille pour être exécuté, que les exécutions sont acceptées quand elles se présentent, que les deux côtés
sont cotés, et que la logique d'annulation relève de la politique, documentée. Le moteur est l'endroit où cette preuve est produite~: il consigne, pour chaque ordre, la
cible qu'il servait.

\section{Tutoriel~: un moteur, quatre configurations}\label{tut:hf:the-quoting-engine:tut}

\begin{tutorial}
\textbf{Objectif.} Piloter le moteur à partir d'une politique de cotation sur le marché simulé et mesurer les messages, les exécutions, les courses, les places dans la file et l'edge.
\textbf{État final~:} le tableau ci-dessous.
\begin{tutsteps}
\item \textbf{Le différentiel} (\cref{met:hf:the-quoting-engine:diff}), dans \texttt{firm.quoteengine.QuoteEngine.update}~; ses jumeaux en C++20 et en Rust produisent les
mêmes 1\,154 actions sur une fixture de 3\,000 événements.
\omcode{code/firm/quoteengine/firm_quoteengine.py}{71}{102}{Couvrir les cibles par des ordres au repos, modifier à la baisse, ajouter, annuler ou conserver dans l'hystérésis.}\label{lst:hf:the-quoting-engine:diff}
\item \textbf{Émission sous budget}~: les annulations d'abord, puis les modifications, puis les nouveaux ordres, un jeton par action.
\omcode{code/firm/quoteengine/firm_quoteengine.py}{103}{124}{Le seau à jetons décide de ce qui est envoyé~; le reste est recalculé à la mise à jour suivante.}\label{lst:hf:the-quoting-engine:issue}
\item \textbf{Le coteur.} \texttt{hf\_engine.EngineQuoter} vise les meilleurs prix des autres (et les prix qui se trouvent derrière eux pour une échelle), un lot par
niveau, et traduit les actions du moteur en ordres~; sur \texttt{firm.tape}, qui n'a pas de message de modification, une modification devient une annulation suivie d'un
nouvel ordre.
\item \textbf{Comparer} quatre configurations sur quatre sessions de vingt minutes avec \texttt{compare()}.
\end{tutsteps}
\textbf{Pour aller plus loin.} Faire tourner le même coteur sur \texttt{firm.exchsim}, dont le message de remplacement conserve la priorité lors d'une réduction de taille~; ajouter une
deuxième stratégie sur le même instrument et compter les auto-appariements empêchés.
\end{tutorial}

\begin{center}
\footnotesize
\setlength{\tabcolsep}{4pt}%
\begin{tabular}{@{}>{\raggedright\arraybackslash}p{3.6cm}>{\centering\arraybackslash}p{1.7cm}>{\centering\arraybackslash}p{1.7cm}>{\centering\arraybackslash}p{1.7cm}>{\centering\arraybackslash}p{1.7cm}>{\centering\arraybackslash}p{1.7cm}>{\centering\arraybackslash}p{1.7cm}@{}}
\toprule
configuration & messages & exécutions & courses & mark-out à 10~s & edge & âge médian\\
& par heure & par heure & par heure & (ticks par action) & (\$ par heure) & d'une exécution (s)\\
\midrule
un niveau & 758 & 250 & 19.5 & 0.141 & 35.25 & 9.6\\
trois niveaux & 980 & 282 & 20.3 & 0.144 & 40.50 & 17.3\\
trois niveaux, hystérésis de deux ticks & 806 & 364 & 1.5 & $-0.064$ & $-23.25$ & 9.6\\
trois niveaux, 2 messages par seconde & 884 & 282 & 9.8 & 0.154 & 43.50 & 17.7\\
\bottomrule
\end{tabular}
\end{center}

Quatre enseignements (\cref{fig:hf:the-quoting-engine:configs}). L'échelle envoie 29~\% de messages de plus et rapporte 15~\% d'edge de plus par heure, parce que ses
exécutions proviennent d'ordres qui ont attendu dans la file (un âge médian de 17 secondes contre 10). L'hystérésis économise des messages, mais elle maintient des ordres à un tick
de l'endroit où la politique les veut, y compris un cours acheteur situé un tick au-dessus d'un meilleur cours acheteur qui a baissé~: cet ordre reste seul à un prix périmé et se fait
sélectionner, et le mark-out devient négatif~; dans un carnet à un tick, l'hystérésis ne doit jamais conserver les ordres devant la cible, seulement derrière. Le limiteur de deux messages par seconde fait tomber 1\,396 actions dérivées par heure
sans rien perdre~: les actions écartées avaient été recalculées puis envoyées lorsqu'un jeton était disponible, et le budget a réduit les courses de moitié. Les messages par exécution,
de 3.0 à 3.5, sont ce que compterait le ratio ordres sur transactions d'une plateforme.

\begin{omfigure}
\begin{tikzpicture}
\begin{axis}[width=11cm,height=5.4cm,ybar,bar width=10pt,ymin=0,ymax=1100,xmin=-0.5,xmax=3.5,xtick={0,1,2,3},
  xticklabels={un niveau,trois niveaux,hystérésis,limiteur},ylabel={\small par heure},tick label style={font=\footnotesize},
  axis y line*=left,area legend,table/col sep=comma,legend style={font=\scriptsize,at={(0.5,-0.22)},anchor=north},legend columns=3]
\addplot[fill=omThm!60,draw=omThm] table[x=i,y=messages]{figdata/hft/10-the-quoting-engine/configs.csv};
\addplot[fill=gray!40,draw=gray] table[x=i,y=fills]{figdata/hft/10-the-quoting-engine/configs.csv};
\addlegendentry{messages envoyés}\addlegendentry{exécutions}
\addlegendimage{line legend,omDef,only marks,mark=*}\addlegendentry{edge, \$ par heure (à droite)}
\end{axis}
\begin{axis}[width=11cm,height=5.4cm,ymin=-40,ymax=60,xmin=-0.5,xmax=3.5,axis y line*=right,axis x line=none,
  ylabel={\small edge (\$ par heure)},tick label style={font=\footnotesize},table/col sep=comma]
\addplot[omDef,only marks,mark=*,mark size=2pt] table[x=i,y=edge]{figdata/hft/10-the-quoting-engine/configs.csv};
\draw[omDef!50,dotted] (axis cs:-0.5,0) -- (axis cs:3.5,0);
\end{axis}
\end{tikzpicture}

\omcaption{Messages et exécutions par heure (barres, à gauche) et edge à dix secondes en dollars par heure (points, à droite) pour les quatre configurations du
tutoriel~: un niveau, une échelle à trois niveaux, l'échelle avec une hystérésis de deux ticks, et l'échelle sous un limiteur de deux messages par seconde avec une
rafale de cinq~; un lot par niveau, quatre sessions de vingt minutes. Données~: \texttt{hf\_engine.compare}.}\label{fig:hf:the-quoting-engine:configs}
\end{omfigure}

\section{Construction~: le moteur de cotation}\label{bld:hf:the-quoting-engine:quoteengine}

\begin{build}
\bfield{Rôle} Transformer les cibles de cotation en un minimum de messages selon les règles de priorité d'une plateforme et un budget de messages, et suivre l'état de chaque ordre.
\bfield{Interface} \texttt{QuoteEngine(min\_move, min\_size, rate, burst)}~: \texttt{update(t, side, targets) -> actions}
(\texttt{('new', oid, side, price, qty)}, \texttt{('amend', oid, qty)}, \texttt{('cancel', oid)})~; \texttt{ack(oid)}, \texttt{fill(oid, qty)}, \texttt{orders}, \texttt{stats}
(new, amend, cancel, dropped, fills, races)~; \texttt{TokenBucket(rate, burst).take(t)}. C++20~: \texttt{cpp/firm\_quoteengine.hpp}~; Rust~:
\texttt{rust/}.
\bfield{Règles} Ne réduire qu'à la baisse~; les ordres en transit ne sont pas modifiés~; annulations avant modifications avant nouveaux ordres~; un seau vide écarte les actions,
sans jamais les mettre en file~; une exécution pendant une annulation en attente est comptée comme une course.
\bfield{Tests d'acceptation} \texttt{code/firm/quoteengine/tests/}~: aucun message lorsque le carnet correspond déjà à la cible~; un déplacement de prix annule
un ordre et en ajoute un~; une réduction de taille modifie, une augmentation ajoute~; l'hystérésis garde un ordre à moins d'un tick~; le limiteur écarte et la mise à jour suivante
recalcule~; une course annulation--exécution est comptée~; les moteurs en Python, C++20 et Rust produisent le même journal d'actions sur la fixture.
\bfield{Pour aller plus loin} Remplacer les messages par la règle de priorité de la plateforme sur \texttt{firm.exchsim}~; des groupes d'auto-appariement par stratégie~; un budget de ratio ordres sur transactions par
instrument et par session.
\end{build}

\begin{omsources}
\item Commission Delegated Regulation (EU) 2017/566 sur le ratio entre les ordres non exécutés et les transactions.
\item CME Group, CME Globex Messaging Efficiency Program, dépôt de règle auprès de la CFTC, novembre 2019.
\item \emph{United States v. Coscia}, 866 F.3d 782 (7th Cir. 2017).
\end{omsources}

\section{Exercices}

\begin{exercise}[$\star$]\label{exo:hf:the-quoting-engine:1}
Le moteur laisse au repos 300 actions à 100 et la cible devient 100 actions à 100. Qu'envoie-t-il, et l'ordre conserve-t-il sa place~?
\end{exercise}

\begin{exercise}[$\star$]\label{exo:hf:the-quoting-engine:2}
Un membre a envoyé 12\,000 ordres et réalisé 400 transactions sur un instrument au cours d'une session. Quel est son ratio d'ordres non exécutés
sur transactions, en nombre~?
\end{exercise}

\begin{exercise}[$\star$]\label{exo:hf:the-quoting-engine:3}
Un seau à jetons se recharge à raison de 2 jetons par seconde avec une profondeur de 5. Il est plein en $t=0$ et cinq actions sont envoyées d'un coup. Combien peuvent être envoyées à
$t=1.2$~?
\end{exercise}

\begin{exercise}[$\star\star$]\label{exo:hf:the-quoting-engine:4}
Pourquoi un moteur limité devrait-il écarter les actions plutôt que les mettre en file~?
\end{exercise}

\begin{exercise}[$\star\star$]\label{exo:hf:the-quoting-engine:5}
Pourquoi l'hystérésis de deux ticks a-t-elle rendu le mark-out négatif, et comment le corriger sans perdre son économie de messages~?
\end{exercise}

\begin{exercise}[$\star\star$]\label{exo:hf:the-quoting-engine:6}
Qu'est-ce qui rend une \omterm{def:hf:the-quoting-engine:ladder}{échelle de cotations} défendable face à une accusation de layering~?
\end{exercise}

\begin{exercise}[$\star\star\star$]\label{exo:hf:the-quoting-engine:7}
\emph{Programmation.} Rejouer la fixture du composant dans un moteur avec \texttt{rate=5} et \texttt{burst=2}
(\texttt{hf\_engine.fixture\_stats}), rapporter les comptes et expliquer pourquoi ils ne disent rien de ce que ferait un limiteur plus strict.
\end{exercise}

\begin{exercise}[$\star\star\star$]\label{exo:hf:the-quoting-engine:8}
\emph{Trouver la faille.} «~Pour économiser des messages, nous augmentons la taille d'un ordre par une modification plutôt que d'envoyer un second ordre.~»
\end{exercise}

\section{Problème~: Dix mille changements d'avis}

\begin{problem}[{Problème du week-end --- dix mille changements d'avis}]\label{pb:hf:the-quoting-engine:1}
Le moteur d'un teneur de marché doit maintenir à jour une échelle à trois niveaux sur une action très active à un tick, sans dépasser son budget de messages ni le
ratio ordres sur transactions de la plateforme.

\medskip\noindent\textbf{Partie I --- Le moteur.}
\begin{enumerate}
\item Énoncer le diff de la \cref{met:hf:the-quoting-engine:diff} et justifier l'ordre dans lequel les actions sont envoyées.
\item Définir une \omterm{def:hf:the-quoting-engine:ladder}{échelle de cotations} et dire à quoi servent ses niveaux profonds dans un carnet à un tick.
\item Définir une \omterm{def:hf:the-quoting-engine:replace}{annulation-remplacement} et la règle de priorité, et expliquer pourquoi le moteur ne modifie jamais à la hausse.
\item Qu'est-ce qu'une course entre annulation et exécution, et à quelle fréquence s'est-elle produite ici~?
\end{enumerate}

\medskip\noindent\textbf{Partie II --- Limites.}
\begin{enumerate}[resume]
\item Définir un \omterm{def:hf:the-quoting-engine:throttle}{limiteur de messages} et décrire le seau à jetons.
\item Résumer les deux limites de messages de l'encadré daté.
\item Qu'est-ce que la prévention de l'auto-appariement (self-trade prevention) et pourquoi une société multistratégies en a-t-elle besoin~?
\item Où se situe la frontière entre une échelle et du layering, et qu'a dit le tribunal dans l'affaire Coscia au sujet des taux d'annulation~?
\end{enumerate}

\medskip\noindent\textbf{Partie III --- Mesures.}
\begin{enumerate}[resume]
\item Donner les messages, les exécutions et l'edge pour les cotateurs à un niveau et à trois niveaux.
\item Pourquoi les exécutions de l'échelle proviennent-elles d'ordres plus anciens~?
\item Qu'a fait l'hystérésis de deux ticks, et pourquoi~?
\item Qu'a rejeté le limiteur, et quel en a été le coût~?
\end{enumerate}

\medskip\noindent\textbf{Partie IV --- Le verdict.}
\begin{enumerate}[resume]
\item Énoncer le \emph{résultat nommé}~: les messages par heure et l'âge médian d'une exécution selon les configurations à un niveau, à trois niveaux et limitée,
face à l'edge que chacune rapporte.
\item Quelle configuration faut-il faire tourner, et avec quelle modification de l'hystérésis~?
\item Que signifieraient les chiffres de messages par exécution pour le ratio ordres sur transactions d'une plateforme~?
\item Comment les résultats changeraient-ils sur une plateforme dont le remplacement conserve la priorité lors des baisses de taille~?
\item Pourquoi les moteurs Python, C++ et Rust doivent-ils concorder action par action~?
\item Que consigne le moteur pour défendre la cotation de la société~?
\item Comment une deuxième stratégie sur le même titre modifierait-elle le moteur~?
\item En une phrase~: qu'optimise un moteur de cotation~?
\end{enumerate}
\end{problem}

\section{Questions d'entretien}

\begin{interviewq}[$\star$ \iqroles{développeur}]\label{iq:hf:the-quoting-engine:1}
Votre moteur envoie une annulation et reçoit aussitôt une exécution pour le même ordre. Quels états doit-il gérer, et que dit-il à la stratégie~?
\end{interviewq}

\begin{interviewq}[$\star\star$ \iqroles{développeur}]\label{iq:hf:the-quoting-engine:2}
Concevez la structure de données des ordres en cours pour que le diff avec une échelle à dix niveaux s'exécute en temps constant par niveau.
\end{interviewq}

\begin{interviewq}[$\star\star$ \iqroles{trader}]\label{iq:hf:the-quoting-engine:3}
Votre plateforme vous impose une limite de 50 messages par seconde, et le marché bouge vite. Quels messages envoyez-vous en premier~?
\end{interviewq}

\begin{interviewq}[$\star\star$ \iqroles{risque}]\label{iq:hf:the-quoting-engine:4}
Un régulateur vous demande pourquoi votre société annule 97~\% de ses ordres. Que lui montrez-vous~?
\end{interviewq}

\begin{interviewq}[$\star\star$ \iqroles{chercheur}]\label{iq:hf:the-quoting-engine:5}
Comment mesureriez-vous la valeur d'une place dans la file que votre échelle préserve lorsque le prix bouge~?
\end{interviewq}

\begin{interviewq}[$\star\star\star$ \iqroles{développeur}]\label{iq:hf:the-quoting-engine:6}
Démontrez qu'un seau à jetons de taux $r$ et de profondeur $b$ ne laisse jamais passer plus de $b+r\,T$ messages sur une fenêtre de longueur $T$.
\end{interviewq}
